Predicting a language model's full-benchmark score from its results on a few items, using per-item results of other models, could cut evaluation cost, but published methods were compared under different protocols.
We evaluate 13 binary-score benchmarks (Open LLM Leaderboard v2, HELM) with one protocol: equal budgets (10--200 items), source/target splits that are random, extrapolating (the strongest 20\% of models as targets) or whole-organisation, source-only selection and tuning, duplicate models removed, and model-level bootstrap intervals with paired comparisons.
We re-implement existing and simple methods and add \mass{}, a model-assisted sampling estimator (stratified probability-proportional-to-size item design, latent-trait predictor, design-weighted correction on the sampled items).
Against the random-subset mean at 50--100 items, \mass{}-Ens lowers RMSE by roughly 21--27\% (all 13 benchmarks; effective-sample-size gain about 1.6--1.8). It is not distinguishable from the best baseline chosen by cross-validation on source models, nor from ridge regression on a random subset when targets resemble the sources, and it does not improve the ranking of target models. Under extrapolation, predictors fitted on weaker models degrade (ridge: RMSE 1.4--1.7 times the random mean), and the difference correction is what prevents this. Its reported intervals under-cover at 10--20 items.
The evidence supports ``effective under conditions'': useful for estimating scores, not for ranking the strongest models.
